\chapter{Conclusion}
\label{ch:conclusion}

This thesis set out to address a specific, concrete problem: a reader interested in AI developments faces content scattered across a dozen heterogeneous source types (research preprints, code releases, video talks, government filings, funding announcements) with no single system that ingests, understands, and personalizes across all of them at once, and AI Compass was designed, built, and evaluated as a working answer to that problem.

The design principle that runs through every system chapter is a deliberate one: generative AI is confined to a small number of well-scoped, independently verifiable roles, and kept out of the ranking and retrieval hot path entirely. A single structured enrichment call per content item, a citation-checked RAG assistant, and a map-reduce video-chaptering pass are the only places a language model is invoked anywhere in the system, and each is wrapped in a mechanical check (validated output fields, server-side citation resolution, transcript-aligned chunk boundaries) that constrains what the model's output is allowed to assert rather than trusting it outright. The ranking core that decides what a user actually sees, and the similarity search underneath both the recommender and the assistant, are fully deterministic and contain no model call at all. This is not an incidental implementation detail; it is the property that makes the rest of the thesis's claims (explainability, groundedness, reproducibility of the non-exploratory ranking output) checkable rather than asserted.

The thesis's contributions follow directly from this principle and are not restated here beyond naming them: a source-adapter ingestion architecture with a single-call, schema-validated enrichment step and hand-rolled near-duplicate clustering (Chapter~5); a deterministic, fully explainable personalized ranking formula with MMR-based diversification and a disclosed exploration mechanism (Chapter~6); a citation-validated, access-controlled retrieval-augmented conversational assistant (Chapter~7); an offline, CPU-only deep-media transcription and chaptering pipeline with a measured throughput figure (Chapter~8); and a deployed three-codebase implementation with an explicitly documented infrastructure evolution (Chapter~9).

The evaluation in Chapter~10 must be stated with the same honesty it was written with. The offline NDCG@10/MAP harness was shown to be correctly implemented (it computes both metrics exactly as textbook-defined, handles the empty-relevant-set case without error, and correctly distinguishes an empty evaluation window from a genuinely poor ranking) and running it surfaced and fixed a real methodological flaw in how the held-out window was anchored. But at the data scale available at evaluation time, every measured NDCG@10 and MAP value was $0.0$, for the structural reason documented in Chapter~10: the corpus of temporally-aligned interaction data needed to produce a meaningful score simply does not yet exist. This is a finding about the current scale of usage, not a failure of the ranking algorithm, and the thesis has been deliberately explicit about the difference throughout rather than allowing one to be mistaken for the other.

What remains is, correspondingly, a matter of scale and accumulated usage rather than of unresolved design. Chapter~11's future work sets out the concrete, bounded path forward: re-running the same evaluation harness once real usage accumulates, closing the ANN-index and re-ingestion gaps identified in that chapter, putting the entity-affinity signal to use in the ranking formula that already computes it, and revisiting collaborative filtering, per-persona summaries, and region- or language-aware filtering once the user base and content pool grow enough to make each of them meaningful. None of these require rethinking what AI Compass is; they extend a system whose core architectural bet (deterministic where it must be verifiable, generative only where its output can be checked) was the thesis's central contribution and remains its starting point for everything that comes next.
